\subsection{Use Case}

In line with \cite{Raji2020_Accountability} and \cite{poretschkin2021_IAIS_Leitfaden} the first step in our framework is a summary of the use case. This offers the auditing authority a brief overview about the task at hand and the AI application. Furthermore, the document can be used during the AI development phase to ensure compliance as well as a trustworthy and robust AI system. A general understanding of the use case is necessary to perform a proper risk assessment afterwards. Therefore, a use case summary should be formulated that includes a description of the status before applying an AI solution, a description of relevant sub-components as well as a summary of the proposed AI solution. This approach is especially useful for the documentation of high-risk applications as defined in the AI Act \citep{AiAct}. A general summary should also include a contact person and information about involved groups inside an organisation to include them in relevant parts during the development phase. New in our framework is the justification of why a non-AI approach is not sufficient in the context of the application. Although this requirement is not explicitly stated in the AI Act, it emerged in discussions with several certification bodies. For developers this means that a simpler model (or maybe non-AI) solution should be prioritized whenever there is no sound justification for the use of a more complex one. Finally, two new aspects in our approach are (i) reviewing and documentation of prior incidents in similar use cases and (ii) the distinction whether a general purpose AI system or a foundation model is expected to be used or developed. 

\subsubsection{Risk} 
AI can only be deployed if the associated risks are kept on an acceptable level, which is comparable to already existing products or services. To guarantee this, a proper risk assessment must be performed. Due to the nature of AI systems, e.g., the training on noisy and changing data or the black-box characteristics, classical methods for risk assessment cannot be applied in most cases \citep{Aria2020_VS, Siebert2022_Quality}. In literature, there are several schemes for qualitative and quantitative risk assessment. In recent years, the AI community proposed several ways to categorize risks in different dimensions \citep{poretschkin2021_IAIS_Leitfaden, Ashmore2021_Assuring, Piorkowski2022_Risk}. \pagebreak 
One approach, accepted by many scientists and organisations, is the division according to the High-Level Expert Group on AI \citep{EU2019_Ethic}. 
They propose the seven key dimensions of risk: 
\emph{
    \begin{enumerate}
        \item human agency/oversight, 
        \item  technical robustness/safety, 
        \item  privacy/data governance, 
        \item  transparency, 
        \item  fairness, 
        \item  environmental/societal well-being, and 
        \item  accountability. 
    \end{enumerate}
}
As described in Section \ref{sec:other_activities}, these dimensions are well-aligned to other works. Based on every dimension one can derive possible risks for one specific use case and AI application. Here, the risk consists of the probability of an undesirable event and the impact such an event has. In recent years, many more methods have been developed to accurately quantify such probabilities, for example, in the field of adversarial robustness \citep{Szegedy2013_Adversarial, Murakonda2020_Privacy}, fairness \citep{Barocas2016_BigDD, Bellamy2018_360, bird2020_fairlearn}, or privacy \citep{Shokri2016_privacy, Fredrikson2015_Inversion}. In the context of (functional) safety, a quantitative assessment of the risk is almost always mandatory, especially in safety critical applications such as robotics, manufacturing, or autonomous driving \citep{El-Shamouty2022_GLIR,Ashmore2021_Assuring, Salay2017_26262}. There are also more qualitative approaches to classify the risk corresponding to one dimension. The \ac{EU} separates between forbidden applications, high risk applications and uncritical applications based on the application domain \citep{AiAct}. %TÜV Austria proposes a scale of four criticality levels as mentioned in section \ref{subsubsec:tuv} \citep{winter2021_TueV_Austria}.
%\ref{sec:other_activities}
A standardized way of performing risk assessment for AI is given in ISO/IEC 23894 which uses the existing risk management standard ISO 31000 as a reference point. It also offers concrete examples in the context of correct risk management integration during and after the development.


Our proposal of a use case documentation is shown below. In accordance with existing works on model cards \citep{Mitchell2019_ModelCards} and data cards \citep{Pushkarna2022_DataCards} we call it a \emph{Use Case Card}. As in the other cards for almost every aspect there exist one or multiple references to the AI Act. After a use case summary and a brief formulation of the problem and solution, we suggest to perform a risk assessment based on the dimensions mentioned above. The risks has to be stated associated with the corresponding dimension. Note that in this approach first a qualitative risk assessment is performed to create awareness of the critical aspects of the AI system. Later on, in the Sections \ref{subsec:data}--\ref{subsec:monitoring}, the goal is to quantify and minimize those risks using specific methods and other measures.

\newpage

\label{tab:uc}
{\small
    \renewcommand\arraystretch{1.3}
    \begin{longtable}{p{0.15\linewidth}|p{0.35\linewidth}p{0.15\linewidth}p{0.35\linewidth}} 
    \toprule
    \multicolumn{4}{c}{\bf{Use Case Card}}\\
    \toprule
    \bf{Step}        & \bf{Requirement}      & \bf{AI Act}& \bf{References} \\
    \hline
    \hline
    \multirow{12}{=}{General}
                        & Name contact person & Art. 16, Art. 24 & \\\cline{2-4}
                        & List groups of people involved & & \\\cline{2-4}
                        & Summarize the use case shortly & Art. 11 & \cite{poretschkin2021_IAIS_Leitfaden}\\\cline{2-4}
                        & Describe the status quo & & \\\cline{2-4}
                        & Describe the planed interaction of sub-components (e.g., different software modules or hardware and software) & Art. 18 & \\\cline{2-4}
                        & Provide a short solution summary & & \cite{Mitchell2019_ModelCards}\\\cline{2-4}      
                        & Review and document past incidents in similar use cases & & \href{https://incidentdatabase.ai/}{AI Incident Database}\\\cline{2-4}
                        & Review available tools & & \href{https://oecd.ai/en/catalogue/overview}{Catalogue of Tools \& Metrics for Trustworthy AI}\\

    \specialrule{.075em}{0em}{0em}
    \multirow{5}{=}{Problem definition}
                        & Clearly describe the learning problem & Art. 11 & \\\cline{2-4}
                        & Describe the disadvantages of the current approach & & \\\cline{2-4}
                        & Argue why classical approaches are not sufficient & & \\
                        
    \specialrule{.075em}{0em}{0em}
    \multirow{7}{=}{Solution approach}
                        & Describe the integration into the current workflow & Art. 11 & \\\cline{2-4}
                        & Formulate the learning problem & & \\\cline{2-4}
                        & Formulate the KPIs for go-live & & \\\cline{2-4}
                        & Document if a foundation model will be used or developed & Art. 3, Art 28 b) & \\\cline{2-4}
                        & Provide a short data description & Art. 10 & 
                        \cite{Gebru2021_DataSheets}\\
                        
    \specialrule{.075em}{0em}{0em}
    \cellcolor[HTML]{d3d3d3}& \cellcolor[HTML]{d3d3d3}\textbf{Risk Assessment} &\cellcolor[HTML]{d3d3d3} Art. 9, Art. 19 &\cellcolor[HTML]{d3d3d3} \\
    \hline
    \hline
    \multirow{2}{=}{Risk Class}          
                        & Categorize the application into a risk class according to the AI Act & Art. 5, Art. 6 & \href{https://blog.burges-salmon.com/post/102i9my/navigating-the-eu-ai-act-flowchart}{Navigating AI Act Flow Chart} \href{https://appliedaiinitiative.notion.site/Risk-Classification-Database-2b58830bb7f54c9d8c869d37bdb27709}{Risk-Classification DB}\\
                        
    \specialrule{.075em}{0em}{0em}
    \multirow{2}{=}{Human agency and oversight} 
                        & Rate the level of autonomy & Art. 14 \\
                                                &                   &\\
    \specialrule{.075em}{0em}{0em}
    \multirow{8}{=}{Technical robustness and safety}
                        & Evaluate the danger to life and health & Art. 5 (1) a), Art. 14 & \cite{winter2021_TueV_Austria} \\\cline{2-4}
                        & Identify possibilities of non-compliance & Art. 15, Art. 19 & \cite{johnson2020understanding}\\\cline{2-4}
                        & Identify customer relevant malfunction & & \\\cline{2-4}
                        \pagebreak
                        \cline{2-4}
                        & Identify internal malfunction & & \\\cline{2-4}
                        & Evaluate cybersecurity risks & Art. 15, Art. 42 & \\
                        
    \specialrule{.075em}{0em}{0em}
    \multirow{5}{=}{Privacy and data governance}
                        & List risks connected with customer data & Art. 10 & \cite{Cristofaro2020_Privacy}\\\cline{2-4}
                        & List risks connected with employee data & & \\\cline{2-4}
                        & List risks connected with company data & & \\
                        
    \specialrule{.075em}{0em}{0em}
    \multirow{4}{=}{Transparency}
                        & Evaluate effects of incomprehensible decisions or the use of a black box model & Art. 13 & \cite{larsson2020transparency, ehsan2021expanding, felzmann2019transparency, rudin2019stop, castelvecchi2016can} \\\cline{2-4}
                        
    \specialrule{.075em}{0em}{0em}
    \multirow{6}{=}{Diversity, non-discrimination and fairness}
                        & Check for possible manipulation of groups of people & & \cite{ISO12791_Fairness, ashton2022problem, botha2020fake}  \\\cline{2-4}         
                        & Check for discrimination of groups of people regarding sensitive attributes & Art. 5 (1) b) and c) &  Section\ref{app:protected_attrs}, \cite{ISO12791_Fairness, Handbook2018_NonDiscrimination, Oneto_2020}\\\cline{2-4}
                        
    \specialrule{.075em}{0em}{0em}
    \multirow{9}{=}{Societal and environmental well being}
                        & List possible dangers to the environment & & \cite{wu2022sustainable} \\\cline{2-4}
                        & Consider ethical aspects & & \cite{dubber2020oxford, muller2020ethics, bostrom2018ethics, hagendorff2020ethics, EU2019_Ethic} \\\cline{2-4}
                        & Evaluate effects on corporate actions & & \cite{glikson2020human}  \\\cline{2-4}
                        & Identify impact on the staff & & \cite{malik2022impact} \\
                        
    \specialrule{.075em}{0em}{0em}
    \multirow{3}{=}{Accountability}
                        & Estimate the financial damage on failure & Art.17 & \cite{gualdi2021artificial} \\\cline{2-4}
                        & Estimate the image damage on failure & & \\
                        
    \specialrule{.075em}{0em}{0em}
    \multirow{3}{=}{Norms}
                        & List relevant norms in the context of the application (e.g., automotive safety norms in an automotive use-case) & Art. 9 (3) & \cite{gasser2020role} \href{https://aistandardshub.org/ai-standards-search/}{AI Standards Hub}\\
                        
    \specialrule{.075em}{0em}{0em}
    \multirow{2}{=}{Involvement of Individuals}    
                        & If feasible, describe involvement of affected individuals  & \\
    \bottomrule
\end{longtable}
}
